\documentclass[conference]{IEEEtran}

\usepackage{cite}
\usepackage{amsmath,amssymb,amsfonts}
\usepackage{graphicx}

\begin{document}

\title{Paper Title}

\author{
\IEEEauthorblockN{Author Name}
\IEEEauthorblockA{
Department, Institution\\
City, Country\\
Email
}
}

\maketitle

\begin{abstract}
% TODO: Abstract.
\end{abstract}

\begin{IEEEkeywords}
% TODO: Keywords.
\end{IEEEkeywords}

\section{Introduction}
\label{sec:introduction}
Global routing determines how wires connnect each pin of a net in physical design. 
The parasitic resistance and capacitance (RC) in each net
could vary drastically across different routing decisions.
These RC further affect the wire delay and the load capacitance seen by the cell.
One bad routing could also increase delay on critical paths and cause timing violations on 
its downstream endpoints. 
However, the routing resource is limited and we cannot always assign 
the best RC for each net. 
Therefore, it is essential for timing driven global routing to identify
which nets are timing critical and how many resources should be allocated to them.

Previous works have proposed various methods to solve this problem. 
GAP-LA \cite{gap-la} first detemines the 2D topology of the routing tree and then 
uses a look ahead function to eastimate the timing cost for the 3D layer assignment.
They define net criticality based on the number of worst timing paths traversing the net and 
minimum pin slack.
Zhou et al. \cite{zhou2026aspdac} balance the timing cost and congestion cost in 
the dynamic programming process of the 3D DAG routing. 
They classify a net as critical if any of its pins has negative slack.
Xiao el at. \cite{xiao2026ispd} 
further apply the log-sum-exponential (LSE) function and derive timing gradient 
to identify the timing critical nets and give each routing segment 
a loss function driven weight. 
However, all these methods focus on optimization under elmore delay and lumped capacitance model. 
They do not explicitly consider cell delay and slew changes caused by different routing choices.  
Moreover, their net criticality and timing cost are evaluated based on one single STA result, neglecting 
the timing changes after some nets have been routed.

In signoff STA, the cell delay is evaluated using effective-load model, such as 
Dartu-Menezes-Pileggi (DMP) \cite{dmp} model. They first reduce the parasitic RC 
using the Pi model \cite{PImodel} to a resistance, a near end capacitance, and a 
far end capacitance. 
Then they compute the effective load based
on effective current and voltage waveform eqautions. 
The effective load can be samller than the lumped capacitance because of the 
far end capacitance
can be shielded by the resistance, 
meaning that larger net resistance might reduce the cell delay.
Utilizing the shielding effect, we could relieve the resouce competence for smaller RC during routing 
and reduce the congestion. 

Besides the timing model, the net timing criticality determines the order how nets are routed and plays an important role in
timing driven global routing.
Previous slack based
criticality only shows where the timing violation net is, without considering how much 
negative slack could be reduced if more routing resources are allocated to the net. 
Nevertheless, slacks of nets change after routing each single net and critical paths may become less important. 
% 配图：shielding effect的真实例子。

Based on above observations, we propose a timing gradient under DMP effective load model,
and we dynamically update the net criticality during the routing.   
The main contributions are as follows:
\begin{itemize}
    \item We propose a timing gradient under the DMP effective l
    oad model.
    \item We dynamically update the net criticality based on estimated changes in negative slack.
    \item We demonstrate the effectiveness of our approach through experiments.
\end{itemize}


\section{Related Work}
\label{sec:related-work}
% TODO: Related work.

\section{Methodology}
\label{sec:methodology}
% TODO: Problem formulation and proposed method.

\section{Experimental Results}
\label{sec:experiments}
% TODO: Experimental setup, baselines, results, and analysis.

\section{Conclusion}
\label{sec:conclusion}
% TODO: Conclusions and future work.

% \section*{Acknowledgment}

\section*{References}
% TODO: Replace this heading with a bibliography when references are added.

\end{document}
